\begin{proposition}[Identical global ranks, different arithmetic coranks]
\label{prop:B-paired-corank}
\status{Refuted rank-only determination; explicit corrected coranks proved}
Put \(S=\Spec\Z[1/(2\cdot11\cdot37\cdot389)]\), and let
\[
\begin{aligned}
 E_{11}&:y^2+y=x^3-x^2-10x-20,\\
 E_{37}&:y^2+y=x^3-x,\\
 E_{389}&:y^2+y=x^3+x^2-2x.
\end{aligned}
\]
The abelian schemes
\(A=E_{389}\times E_{11}\) and \(B=E_{37}\times E_{37}\)
over the same base \(S\) have
\[
\begin{gathered}
 \dim A=\dim B=2,\qquad
 \operatorname{rank}A(\Q)=\operatorname{rank}B(\Q)=2,\\
 \Sha(A/\Q)[2^\infty]=\Sha(B/\Q)[2^\infty]=0.
\end{gathered}
\]
Nevertheless their arithmetic cohomology groups have distinct
coranks:
\[
 \operatorname{corank}_{\Z_2}H^1(S,A)[2^\infty]=1,
 \qquad
 \operatorname{corank}_{\Z_2}H^1(S,B)[2^\infty]=0.
\]
Thus even fixing the number field, coefficient prime, base,
dimension, Mordell--Weil rank and entire primary Sha group does
not determine the requested arithmetic corank.  The local closure
contribution \(t\) in Theorem~\ref{thm:B-Sha-corank} is essential.
\end{proposition}
\begin{proof}
The complete algebraic calculation is
\texttt{part\_b/code/paired\_corank.gp}, with full output in
\texttt{part\_b/output/paired\_corank.txt}.  It uses PARI~2.19.1
and no analytic-rank or BSD calculation.  Global reduction gives
conductors \(11,37,389\) and discriminants \(-11^5,37,389\),
respectively, so all three are elliptic schemes over \(S\).
The exact outputs are
\[
\begin{array}{c|c|c|c}
 E&\operatorname{ellrank}(E)&E(\Q)_{\rm tors}&
             \dim_{\F_2}\operatorname{Sel}_2(E/\Q)\\ \hline
 E_{11}&[0,0,0,[]]&\Z/5&0\\
 E_{37}&[1,1,0,[[-1,0]]]&0&1\\
 E_{389}&[2,2,0,[[-2,0],[4,8]]]&0&2.
\end{array}
\]
For precision, the third \texttt{ellrank} entry is
\(s_0=\dim_{\F_2}\Sha[2]/2\Sha[4]\); it is not in
general \(\dim\Sha[2]\).  The official algorithm description
states that the two-Selmer dimension \(C\), rational two-torsion
dimension \(T\), and \(s_0\) are unconditional, with
\[
 r_2=C-T-s_0,\qquad
 C=T+\operatorname{rank}E(\Q)+\dim_{\F_2}\Sha[2].
\]
Here the lower and upper rank bounds coincide, \(T=s_0=0\),
so these identities give \(\Sha[2]=0\) in every case
\citep[entries \texttt{ellrank} and \texttt{ell2cover}]{BPari2026}.

There is a direct additional certificate of that last assertion.
The program outputs a basis of all everywhere locally soluble
two-covers, with its exact rational maps to the elliptic curve.
The quartics are empty for \(E_{11}\), are
\[
       v^2=u^4-4u^3+6u^2+1
\]
for \(E_{37}\), and are
\[
\begin{aligned}
 v^2&=u^4-4u^3-2u^2+8u+1,\\
 v^2&=u^4-4u^3+10u^2+1
\end{aligned}
\]
for \(E_{389}\).  Every listed basis cover has the rational
point \((u,v)=(0,1)\).  Under the recorded covering maps these
points give \((-1,0)\), \((4,8)\), and \((-2,0)\), respectively.
The program verifies each map by an exact polynomial identity.
Thus every two-Selmer basis class has trivial image in Sha.
The two-descent exact sequence proves \(\Sha[2]=0\) directly,
and a primary torsion group with no elements of order two is zero.
Products preserve both the Mordell--Weil decomposition and Sha,
which proves all the stated global equalities.

It remains to compute the local closure ranks, without a numerical
logarithm.  Each \(E(\Q_2)\) contains an open formal subgroup
isomorphic to \(4\Z_2\), by the formal logarithm on depth
\(r=2>v_2(2)/(2-1)=1\)
\citep[Theorem~IV.6.4(b), p.~132, and Proposition~VII.2.2,
p.~191]{BSilverman2009Local}.
It has finite index because \(E(\Q_2)\) is compact.
A non-torsion rational point remains non-torsion over \(\Q_2\):
an equation \([m]P=0\) after field extension is the same equation
over \(\Q\).  A nonzero multiple in the formal subgroup therefore
has a nonzero logarithm.  The closure of its integer multiples is
its logarithm times \(\Z_2\), of rank one.  Since the ambient
local group itself has rank one, the closures of \(E_{37}(\Q)\)
and \(E_{389}(\Q)\) each have rank one; the finite group
\(E_{11}(\Q)\) has closure rank zero.

The rational and local point groups of a product are products,
and closure commutes with a finite product.  Consequently
\(t(A)=1+0=1\), whereas \(t(B)=1+1=2\).  The product
principal polarizations identify these varieties with their duals.
Because their primary Sha groups vanish, the pro-Selmer groups
in Theorem~\ref{thm:B-Sha-corank} are precisely the completions
of these rational-point groups.  At the other inverted finite
primes the local \(2\)-primary cohomology has corank zero.
Both varieties have \(d=2\) and \(s=0\), so formula (B.20)
gives \(2-1=1\) and \(2-2=0\), as claimed.  The latter group
is finite by the finite-cotype assertion already proved there.
\end{proof}
